\documentclass[preview]{standalone}

\usepackage{amsmath}
\usepackage{amssymb}
\usepackage{stellar}
\usepackage{definitions}

\begin{document}

\id{wundt-e-introspezione}
\genpage

\section{Wundt e l'introspezione}

\begin{snippet}{wundt-psicologia-moderna}
    Con \emph{Wilhelm Wundt} nasce la psicologia moderna: Wundt applica i
    metodi della fisiologia ai processi e ai contenuti della coscienza umana.
    Il suo laboratorio di Lipsia, fondato nel 1879, segna l'avvio della
    psicologia sperimentale come disciplina autonoma.
\end{snippet}

\begin{snippetdefinition}{esperienza-diretta-definition}{Esperienza diretta}
    Per Wundt, l'\emph{esperienza diretta}, o esperienza immediata, è
    l'oggetto di studio della psicologia. Essa riguarda ciò che il soggetto
    vive nel momento in cui si verifica un evento mentale, prima che tale
    esperienza venga trasformata in una spiegazione o interpretazione.
\end{snippetdefinition}

\begin{snippetdefinition}{introspezione-definition}{Introspezione}
    L'\emph{introspezione} è un metodo scientifico usato da Wundt per studiare
    l'esperienza diretta. Richiede il controllo accurato dello stimolo capace
    di produrre l'evento mentale osservato e la stesura di un resoconto subito
    dopo l'osservazione dell'evento.
\end{snippetdefinition}

\begin{snippet}{wundt-linee-guida-introspezione}
    Per rendere scientificamente valida l'introspezione, Wundt richiede che
    l'osservatore, quando possibile:
    \begin{itemize}
        \item possa stabilire quando introdurre il processo osservato;
        \item si trovi in una condizione di sforzo attentivo;
        \item possa ripetere diverse volte la stessa osservazione.
    \end{itemize}
    La condizione sperimentale deve inoltre variare in modo controllato
    l'intensità e la qualità della stimolazione. L'attendibilità del metodo
    dipende dall'abilità dell'osservatore e dalla sua rapidità nel produrre il
    resoconto.
\end{snippet}

\begin{snippetexample}{introspezione-example}{Introspezione}
    In un esperimento introspettivo, lo sperimentatore presenta al partecipante
    uno stimolo controllato, per esempio un suono o un colore, e gli chiede di
    riferire le proprie esperienze immediate. Il resoconto deve riguardare i
    processi percettivi, evitando interpretazioni soggettive. Nel caso di un
    suono, il partecipante può descriverne l'intensità, ma non deve riportare
    valutazioni personali come fastidio, paura o gradimento.
\end{snippetexample}

\begin{snippet}{wundt-limiti-introspezione}
    Secondo Wundt, l'auto-osservazione sperimentale permette di studiare
    processi mentali come appercezione, volontà ed emozioni. Tuttavia, non è
    adatta allo studio dei processi mentali di ordine superiore.
\end{snippet}

\end{document}
